{\large\textbf{Precession of the Earth's axis (10.0 points)}}

\textbf{Introduction}

It has been known since ancient times that the Earth's axis of rotation
precesses. That is, the axis itself rotates around the line perpendicular to
the ecliptic plane, i.e., the plane containing the Earth's orbit around the
Sun. Ancient Greek astronomer \textit{Hipparchus} concluded that the annual
angular displacement of the axis was approximately
45\textquotesingle\textquotesingle{} (seconds of arc), which would imply that
the period of axial precession is around 29000 years. Modern measurements
indicate that the period is approximately 25800 years. In this problem, you are
asked to investigate this phenomenon using Newtonian mechanics.

You may need the following constants:

\begin{itemize}
\item gravitational constant: $G = 6.67 \times 10^{-11}\,\mathrm{Nm^2/kg^2}$
\item average radius of Earth: $R = 6.371 \times 10^{6}\,\mathrm{m}$
\item mass of the Earth: $M_E = 5.972 \times 10^{24}\,\mathrm{kg}$
\item average distance of the Sun from the Earth: $d_{SE} = 1.496 \times 10^{11}\,\mathrm{m}$
\item mass of the Sun: $M_S = 1.989 \times 10^{30}\,\mathrm{kg}$
\item average distance of the Moon from the Earth: $d_{ME} = 3.844 \times 10^{8}\,\mathrm{m}$
\item mass of the Moon: $M_M = 7.348 \times 10^{22}\,\mathrm{kg}$
\item Earth's axial tilt: $\alpha = 23.5^\circ$
\end{itemize}

\textbf{Part A. The shape of the Earth (1.0 points)}

The Sun and the Moon exert nonzero torques on the Earth because of its
non-spherical shape, giving rise to its axial precession. The main reason
behind the Earth's non-spherical shape is the centrifugal force caused by the
Earth's rotation about its axis. The tectonic plates located on the Earth's
surface have deformed over millions of years to minimize stress within them.
Therefore, as an approximation, let us model the Earth as a large liquid
droplet of uniform density whose shape is determined by centrifugal and
gravitational forces. In this model, the Earth's surface is an oblate spheroid
(ellipsoid of revolution) characterized by the polar radius $R_p$ and the
equatorial radius $R_e$ (see \textit{Figure A.1}).

\begin{center}\includegraphics[width=0.28\linewidth]{figures/APhO_2025_Q1_fig1.png}\end{center}

\textit{Figure A.1.} The ellipsoidal shape of the Earth. The polar and
equatorial radii are indicated. $\alpha = 23.5^\circ$ is the angle between the
Earth's axis of rotation and the normal of the ecliptic plane.

The difference between the equatorial and polar radii of the Earth,
$h_{\mathrm{max}} = R_e - R_p$ is much smaller than the average radius
$R = (R_e + R_p)/2$. Up to a dimensionless factor, the value of
$h_{\mathrm{max}}$ can be expressed in terms of the angular speed of the Earth's
rotation $\omega$, its mass $M_E$ and average radius $R$ as
\[
h_{\mathrm{max}} \propto G^{-1} \omega^{\beta} M_E^{\gamma} R^{\delta},
\]
where $G$ is the gravitational constant, and $\beta$, $\gamma$ and $\delta$ are
constant exponents.

\textbf{A.1}\quad Find the values of exponents $\beta$, $\gamma$ and
$\delta$. \hfill 0.8 pt

\textbf{A.2}\quad Calculate the numerical value of $h_{\mathrm{max}}$ assuming
that the dimensionless factor in the relation given above equals 1.
\hfill 0.2 pt

Regardless of whether you were able to find $h_{\mathrm{max}}$ in part
\textbf{A.2.}, use the empirical value $h_{\mathrm{max}} = 21\,\mathrm{km}$ in the
following questions.

\textbf{Part B. The time-averaged gravitational field of the Sun (3.2 points)}

To see why the Sun exerts a nonzero torque (with respect to the center of the
Earth) on our planet, consider \textit{Figure B.1} below. The difference in
distance from the Sun causes the gravitational force $F_1$ to be greater than
its counterpart $F_2$.

\begin{center}\includegraphics[width=0.46\linewidth]{figures/APhO_2025_Q1_fig2.png}\end{center}

\textit{Figure B.1.} Explanation of the nonzero torques exerted by the Sun
(right side of the figure) on the Earth (left side).

The magnitude of this torque acting on the Earth varies continuously during
the year. In the position shown in \textit{Figure B.1}, the torque is maximal,
a quarter of a year later, due to symmetry, the torque becomes zero. After half
a year, it reaches the maximum again, three-quarters of a year later it is zero
once again, and so on. Since the period of axial precession is much larger than
one year, this time-dependent torque can be approximated well by its one-year
average.

To calculate the average torque exerted by the Sun on the Earth, let us
determine first the time average of the gravitational field generated by the
Sun in the vicinity of the Earth. This average can be calculated as the field
of a uniformly dense mass ring, a \textit{Sun ring}, whose mass equals the mass
of the Sun $M_S$ and whose radius equals the average distance between the Sun
and the Earth $d_{SE}$ (see \textit{Figure B.2}).

\begin{center}\includegraphics[width=0.56\linewidth]{figures/APhO_2025_Q1_fig3.png}\end{center}

\textit{Figure B.2.} Time averaging is equivalent to uniformly spreading the
Sun along the circle of radius $d_{SE}$.

Let our cylindrical coordinate system have the origin at the center of the
Earth, and let the $z$ axis be perpendicular to the ecliptic plane (i.e. the
plane of the ring). The axis of rotation of the Earth makes an angle of
$\alpha = 23.5^\circ$ with the $z$ axis.

\textbf{B.1}\quad Find the direction and magnitude of the gravitational field
generated by the Sun ring at a point on the $z$ axis. Write your answer in terms
of $M_S$, $d_{SE}$, and the coordinate $z$. Assume that $|z| \ll d_{SE}$.
\hfill 1.0 pt

\textbf{B.2}\quad Find the direction and magnitude of the gravitational field
generated by the Sun ring at a point in the ecliptic plane whose distance from
the origin is $r$. Assume that $r \ll d_{SE}$. \hfill 2.2 pt

\textbf{Part C. The torque acting on the Earth (2.6 points)}

In this section, you are asked to determine the torque exerted on the Earth due
to the gravitational field obtained in \textbf{Part B}. For simplicity,
consider the Earth as a rigid body with homogeneous mass distribution. Let us
take into account that the rotational ellipsoid can be imagined as if we
removed excess parts from a sphere with the equatorial radius of Earth $R_e$
(see \textit{Figure C.1}).

\begin{center}\includegraphics[width=0.56\linewidth]{figures/APhO_2025_Q1_fig4.png}\end{center}

\textit{Figure C.1.} The ellipsoidal shape of the Earth can be imagined as if
the excess parts were removed from a complete sphere of radius $R_e$.

\textbf{C.1}\quad Find the mass $m$ of one of the two excess regions indicated
in \textit{Figure C.1}. Express your answer in terms of $h_{\mathrm{max}}$, the
mass of the Earth $M_E$, and its polar radius $R_p$. \hfill 0.8 pt

It can be shown that the torque acting on the excess regions is equivalent to
the torque acting on two point masses, each with a mass equal to $2m/5$,
positioned at the endpoints $A$ and $B$ of the polar diameter (see
\textit{Figure C.1}).

\textbf{C.2}\quad Given this idea, find the torque $\tau$ exerted by the Sun
ring on the Earth. Express your answer in terms of $M_E$, $M_S$, $d_{SE}$, $R$
(the average radius), $h_{\mathrm{max}}$ and the angle $\alpha$. You can use
that $h_{\mathrm{max}} \ll R$. \hfill 1.8 pt

\textbf{Part D. Angular speed of the precession of the Earth's axis (2.0 points)}

The Earth's axis of rotation moves \textit{very slowly} around the $z$ axis in
a conical motion. That is, it precesses.

\textbf{D.1}\quad Give an expression for the period $T_1$ of precession of the
Earth's axis. Express your answer in terms of $M_S$, $d_{SE}$, the angular
speed $\omega$ of the Earth's rotation, $h_{\mathrm{max}}$, $R$ and $\alpha$.
\hfill 1.8 pt

\textbf{D.2}\quad Calculate the precession period $T_1$ in years. \hfill 0.2 pt

\textbf{Part E. The effect of the Moon (1.2 points)}

The value obtained in \textbf{Part D} is much larger than the observed value.
The reason for this is that so far we have only considered the torque exerted
by the Sun, and neglected the effect of the Moon. In the following
calculations, assume that the Moon's orbit is in the ecliptic plane, and that
the orbit of the Moon around the Earth is a circle of radius $d_{ME}$. Let us
denote the mass of the Moon by $M_M$ and the period of precession in this
modified model by $T_2$.

\textbf{E.1}\quad By what factor $T_2/T_1$ does the period of precession of the
Earth's axis change if we also take into account the torque exerted by the
Moon? Give your answer in terms of $d_{ME}$, $d_{SE}$, $M_S$ and $M_M$.
\hfill 1.0 pt

\textbf{E.2}\quad By substituting the data, calculate the period of precession
$T_2$ in years. \hfill 0.2 pt
